\documentclass[a4paper,10pt]{article}
\usepackage[utf8]{inputenc}
\usepackage{amsmath}
\usepackage{amsfonts}
\usepackage{amssymb}
\usepackage{mathrsfs}
\usepackage{bm}

\newcommand{\abs}[1]{\lvert{#1}\rvert}            % Small absolute value
\newcommand{\order}{O}
\newcommand{\nup}{n_{\uparrow}}
\newcommand{\nupr}{n_{\uparrow}(\bm{r})}
\newcommand{\ndown}{n_{\downarrow}}
\newcommand{\ndownr}{n_{\downarrow}(\bm{r})}

%opening
\title{HFB densities}
\author{Gabriel Wlaz\l{}owski}

\begin{document}

\maketitle

\begin{abstract}
Basic relations used for computation of various densities.
\end{abstract}

\section{General case}

HFB problem in general case is given by:

\begin{align}\label{eq:hfbspin}
\begin{gathered}
\begin{pmatrix}
h_{\uparrow \uparrow}  & h_{\uparrow \downarrow} & 0 & \Delta \\
h_{\downarrow \uparrow} & h_{\downarrow \downarrow}& -\Delta & 0 \\
0 & -\Delta^* &  -h^*_{\uparrow \uparrow}& -h^*_{\uparrow \downarrow} \\
\Delta^* & 0 & -h^*_{\downarrow \uparrow} & -h^*_{\downarrow \downarrow}
\end{pmatrix} 
\begin{pmatrix}
u_{n\uparrow} \\
u_{n\downarrow} \\
v_{n\uparrow} \\
v_{n\downarrow}
\end{pmatrix}
= E_n
\begin{pmatrix}
u_{n\uparrow} \\
u_{n\downarrow} \\
v_{n\uparrow} \\
v_{n\downarrow}
\end{pmatrix} 
\end{gathered}
\end{align}
and corresponding densities are constructed as follow:
\begin{eqnarray}
n_{\sigma}(\bm{r}) &=& \sum_{|E_n|<E_c}\abs{v_{n,\sigma}(\bm{r})}^2 f_{\beta}(-E_n),\label{eqn:densities0}\\
\tau_{\sigma}(\bm{r}) &=& \sum_{|E_n|<E_c}\abs{\nabla v_{n,\sigma}(\bm{r})}^2 f_{\beta}(-E_n),\\
\nu(\bm{r}) &=& \dfrac{1}{2}\sum_{|E_n|<E_c} (u_{n,\uparrow}(\bm{r})v_{n,\downarrow}^{*}(\bm{r})-u_{n,\downarrow}(\bm{r})v_{n,\uparrow}^{*}(\bm{r}))f_{\beta}(-E_n),\\
\bm{j}_{\sigma}(\bm{r}) &=& \sum_{|E_n|<E_c} \textrm{Im}[v_{n,\sigma}(\bm{r})\nabla v_{n,\sigma}^*(\bm{r})] f_{\beta}(-E_n),
\label{eqn:densities}
\end{eqnarray}
where subscripts $\sigma=\{\uparrow,\downarrow\}$ indicate spin, $E_{n}$ denotes quasi-particle energy and $E_c$ is energy cut-off scale. Fermi distribution function $f_{\beta}(E)=1/(\exp(\beta E)+1)$ is introduced to model temperature $k_B T=1/\beta$ effects. In zero temperature limit $T=0$ Fermi distribution function reduces to  $f_{\beta}(E)=\theta(-E)$, where $\theta(x)$ is Heaviside step function. Note also that $f_{\beta}(E)=1-f_{\beta}(-E)$.\\

\noindent\textit{Relation 1:} If vector $\varphi_{+}=\left( u_{n\uparrow},u_{n\downarrow},v_{n\uparrow},v_{n\downarrow}\right) ^{T}$ is solution of Eq.~(\ref{eq:hfbspin}) with eigenvalue $E_n$, then also vector $\varphi_{-}=(v_{n\uparrow}^*,v_{n\downarrow}^*,u_{n\uparrow}^*,u_{n\downarrow}^*)^{T}$ is solution with eigenvalue $-E_n$.\\

\noindent\textit{Conclusion 1:} It is sufficient to extract only solutions for positive eigenvalues $E_n$. 

\section{Case where $h_{\uparrow \downarrow}=h_{\downarrow \uparrow}=0$}
If $h_{\uparrow \downarrow}=h_{\downarrow \uparrow}=0$ (typically it corresponds to no spin-orbit case) then Eq.~(\ref{eq:hfbspin}) decouples into two sets of equations:

\begin{align}\label{eq:hfbspin1}
\begin{gathered}
\begin{pmatrix}
h_{\uparrow \uparrow}  &  \Delta \\
\Delta^* &  -h^*_{\downarrow \downarrow}
\end{pmatrix} 
\begin{pmatrix}
u_{n\uparrow} \\
v_{n\downarrow}
\end{pmatrix}
= E_n
\begin{pmatrix}
u_{n\uparrow} \\
v_{n\downarrow}
\end{pmatrix} ,
\end{gathered}
\end{align}

\begin{align}\label{eq:hfbspin2}
\begin{gathered}
\begin{pmatrix}
h_{\downarrow \downarrow}& -\Delta \\
-\Delta^* &  -h^*_{\uparrow \uparrow}
\end{pmatrix} 
\begin{pmatrix}
u_{n\downarrow} \\
v_{n\uparrow}
\end{pmatrix}
= E_n
\begin{pmatrix}
u_{n\downarrow} \\
v_{n\uparrow}
\end{pmatrix} .
\end{gathered}
\end{align}

\noindent\textit{Relation 2:} If vector $\varphi_{+}=\left( u_{n\uparrow},v_{n\downarrow}\right)^{T}$ is solution of Eq.~(\ref{eq:hfbspin1}) with eigenvalue $E_n$, then vector $\varphi_{-}=(v_{n\uparrow}^*,u_{n\downarrow}^*)^{T}$ is solution of Eq.~(\ref{eq:hfbspin2}) with eigenvalue $-E_n$.\\

\noindent\textit{Conclusion 2:} Having solutions of Eq.~(\ref{eq:hfbspin1}) one can construct solutions of Eq.~(\ref{eq:hfbspin2}) by using transformation: $u_{n,\uparrow}\rightarrow v_{n,\uparrow}^*$, $v_{n,\downarrow}\rightarrow u_{n,\downarrow}^*$ and $E_n\rightarrow -E_n$. \\
In practice we solve only Eq.~(\ref{eq:hfbspin1}) and we express all densities via $\{u_{n\uparrow},v_{n\downarrow}\}$:
\begin{eqnarray}
n_{\uparrow}(\bm{r}) &=& \sum_{|E_n|<E_c}\abs{u_{n,\uparrow}(\bm{r})}^2 f_{\beta}(E_n),\\
n_{\downarrow}(\bm{r}) &=& \sum_{|E_n|<E_c}\abs{v_{n,\downarrow}(\bm{r})}^2 f_{\beta}(-E_n)\\
\tau_{\uparrow}(\bm{r}) &=& \sum_{|E_n|<E_c}\abs{\nabla u_{n,\uparrow}(\bm{r})}^2 f_{\beta}(E_n),\\
\tau_{\downarrow}(\bm{r}) &=& \sum_{|E_n|<E_c}\abs{\nabla v_{n,\downarrow}(\bm{r})}^2 f_{\beta}(-E_n),\\
\nu(\bm{r}) &=& \dfrac{1}{2}\sum_{|E_n|<E_c} u_{n,\uparrow}(\bm{r})v_{n,\downarrow}^{*}(\bm{r})\left( f_{\beta}(-E_n)-f_{\beta}(E_n)\right) ,\\
\bm{j}_{\uparrow}(\bm{r}) &=& -\sum_{|E_n|<E_c} \textrm{Im}[u_{n,\uparrow}(\bm{r})\nabla u_{n,\uparrow}^*(\bm{r})] f_{\beta}(E_n),\\
\bm{j}_{\downarrow}(\bm{r}) &=& \sum_{|E_n|<E_c} \textrm{Im}[v_{n,\downarrow}(\bm{r})\nabla v_{n,\downarrow}^*(\bm{r})] f_{\beta}(-E_n),\\
S&=&-k_B\sum_{|E_n|<E_c}\left( f_{\beta}(E_n)\ln\ f_{\beta}(E_n) + f_{\beta}(-E_n)\ln f_{\beta}(-E_n)\right).  
\label{eqn:densities1}
\end{eqnarray}

\section{Case where $h_{\uparrow \uparrow}=h_{\downarrow \downarrow}=h$}
If $h_{\uparrow \uparrow}=h_{\downarrow \downarrow}=h$ (typically it corresponds to the spin-balanced case) then Eq.~(\ref{eq:hfbspin1}) takes form, as typically cited in papers (spin indices are dropped):
\begin{align}\label{eq:hfbspin3}
\begin{gathered}
\begin{pmatrix}
h  &  \Delta \\
\Delta^* &  -h^*
\end{pmatrix} 
\begin{pmatrix}
u_{n} \\
v_{n}
\end{pmatrix}
= E_n
\begin{pmatrix}
u_{n} \\
v_{n}
\end{pmatrix} .
\end{gathered}
\end{align}

\noindent\textit{Relation 3:} If vector $\varphi_{+}=\left( u,v\right)^{T}$ is solution of Eq.~(\ref{eq:hfbspin3}) with eigenvalue $E_n$, then vector $\varphi_{-}=(v_{n}^*,-u_{n}^*)^{T}$ is also solution with eigenvalue $-E_n$.\\

\noindent\textit{Conclusion 3:} Having solutions of Eq.~(\ref{eq:hfbspin3}) for positive energies $E_n$ one can construct all solutions of this equation.\\
In practice for spin-balanced case we solve Eq.~(\ref{eq:hfbspin3}) for positive $E_n$ and we express all densities via these states:
\begin{eqnarray}
n_{\uparrow}(\bm{r})=n_{\downarrow}(\bm{r}) &=& \sum_{0<E_n<E_c}\left( \abs{v_{n}(\bm{r})}^2 f_{\beta}(-E_n)+\abs{u_{n}(\bm{r})}^2 f_{\beta}(E_n)\right) \\
\tau_{\uparrow}(\bm{r})=\tau_{\downarrow}(\bm{r}) &=& \sum_{0<E_n<E_c}\left( \abs{\nabla v_{n}(\bm{r})}^2 f_{\beta}(-E_n) + \abs{\nabla u_{n}(\bm{r})}^2 f_{\beta}(E_n)\right) ,\\
\nu(\bm{r}) &=& \sum_{0<E_n<E_c} u_{n}(\bm{r})v_{n}^{*}(\bm{r})\left( f_{\beta}(-E_n)-f_{\beta}(E_n)\right) ,\\
\bm{j}_{\uparrow}(\bm{r})=\bm{j}_{\downarrow}(\bm{r}) &=& \sum_{0<E_n<E_c}\left(  \textrm{Im}[v_{n}(\bm{r})\nabla v_{n}^*(\bm{r})] f_{\beta}(-E_n)\right.\nonumber \\
&& \quad\quad\quad\quad - \left. \textrm{Im}[u_{n}(\bm{r})\nabla u_{n}^*(\bm{r})] f_{\beta}(E_n) \right), \\
S&=&-2k_B\sum_{0<E_n<E_c}\left( f_{\beta}(E_n)\ln\ f_{\beta}(E_n)\right. \nonumber \\
&& \quad\quad\quad\quad\quad\quad + \left. f_{\beta}(-E_n)\ln f_{\beta}(-E_n)\right).
\label{eqn:densities2}
\end{eqnarray}

\end{document}


